\section{无穷大动量态中的部分子量子与大动量展开}
\label{sec:imf}

让我们从部分子形式体系的简要回顾开始。教科书中，PDF 通常由 QCD 中的 LF 关联来定义~\cite{Sterman:1994ce,Collins:2011zzd}。若一个无质量粒子沿 $z$ 方向运动，则光前（LF）由 $t-z=$~常数定义；其他坐标 $t+z$ 与横向空间维度的变化张成一个三维平面。引入两个独立的 LF 四矢量和量纲为一的参数 $\Lambda$，
\begin{eqnarray}
   p^\mu &=& (\Lambda, 0, 0, \Lambda )\ , \nonumber \\
   n^\mu &=& \left(1/(2\Lambda), 0, 0, -1/(2\Lambda)\right)\ ,
\end{eqnarray}
则 $p^2=n^2=0$，且 $p\cdot n=1$。不同的 LF 由沿 $n$ 方向不同的坐标距离 $\lambda$ 定义。

现在考虑质量为 $M$、四动量为
$P^\mu= (P^0,0,0,P^z)
= p^\mu + (M^2/2) n^\mu$ 的态 $|P\rangle$ 中的夸克 PDF，由此可解出 $\Lambda =(P^0+P^z)/2$。
用 $\psi$ 表示完整 QCD 的夸克场，坐标空间的 LF 关联函数为
\begin{equation}
     h(\lambda)= \frac{1}{2P^+}\langle P|\overline \psi(\lambda n)\Gamma W(\lambda n, 0)\psi(0)|P\rangle \ ,
\label{Eq:correlationm}
\end{equation}
其中 $\Gamma$ 是 Dirac 矩阵，$W$ 是直的 Wilson 线规范连接，$\lambda$ 为 LF 距离，其余坐标均取为零。由于 LF 在沿 $z$ 方向的 Lorentz 推进下不变，上述关联函数与剩余动量 $P^z$ 无关。通常取 $P^z=0$。

夸克 PDF 正是上述 LF 关联的傅里叶变换~\cite{Collins:2011zzd},
\begin{equation}
     f(x) =\int^\infty_{-\infty} \frac{d\lambda}{2\pi} e^{-ix\lambda}h(\lambda)  \ .
\end{equation}
这样，尽管部分子是描述 LF 共线模的有效自由度（dof），人们无需借助 EFT 的机器便可研究它们。原因在于：作用于完整 QCD 态 $|P\rangle$ 上的 LF 关联子会自动投影出部分子自由度。另一方面，这些部分子自由度也可在 QCD 拉氏量中被显式分离出来，正如软共线有效理论（SCET）中以 LF 共线场所表示的那样~\cite{Bauer:2000yr,Bauer:2001ct,Bauer:2001yt}。

在传统的部分子形式体系中，关联是含时的；换言之，算符处于 Heisenberg 绘景。因此我们说该形式体系是 Minkowski 式的，由于著名的``符号''问题而难于进行蒙特卡罗模拟。若选取 $\xi^- = (t+z)/\sqrt{2}$ 作为``新时间''坐标并将其积掉，便得到部分子的哈密顿量形式体系，文献中称之为 LF 量子化（LFQ）~\cite{Dirac:1949cp}。尽管已取得很大进展~\cite{Brodsky:1997de}，LFQ 仍是一种极难处理的形式体系。

另一种部分子形式体系可以通过把 Feynman 关于部分子的原始想法移植到场论的语境中获得~\cite{Ji:2014gla,Ji:2020ect}。Feynman 考虑过复合系统的动量分布~\cite{Feynman:1973xc}
$f(k^z, P^z)$，其中 $P^z$ 是质心动量，$k^z$ 是部分子携带的纵向动量，其横动量 $\vec{k}_\perp$ 已被积掉。动量分布对 $P^z$ 的依赖显然是一种相对论性效应：按照 Poincar\'e 对称性，系统的哈密顿量依赖参考系，并在 Lorentz 推进下按
\begin{equation}
  [H, K^i] = i P^i\ ,
\end{equation}
变化，其中 $K^i$ $(i = 1,2,3)$ 是推进算符。因此波函数依赖参考系，导致依赖参考系的动量分布。由于 $K^i$ 依赖于相互作用，这种对参考系的依赖是一个动力学问题，一般需要非微扰求解。

系统动量分布的一个重要特征在于它是一个静态或不含时的量。在 QCD 中，它与如下空间关联相联系，
\begin{equation}
    \tilde h(z,P^z)= \frac{1}{2N}\langle P|\overline \psi(z)\Gamma W(z, 0)\psi(0)|P\rangle \ ,
\label{eq:esme}
\end{equation}
其中 $W(z,0)$ 是类空的直线规范连接，$N$ 是依赖于 Dirac 矩阵的归一化因子。Feynman 进而考虑了无穷大动量极限（假定这样的极限存在），
\begin{equation}
  P^z \to \infty, ~~~~~z\to 0, ~~~~~\lambda = zP^z ~{\rm finite}\ ,
\end{equation}
即与部分子相关的关联函数为
\begin{equation}
  \tilde h(\lambda= zP^z) = \langle P^z=\infty|\overline \psi(z)\Gamma W(z, 0)\psi(0)|P^z=\infty\rangle\ .
\label{Eq:correlatione}
\end{equation}
显然在场论中这是一个非平庸的极限。事实上可以证明，这样的极限只存在于渐近自由理论中，在那里高动量模是微扰的~\cite{Ji:2020ect}。

若忽略该极限的微妙之处，式~(\ref{Eq:correlatione}) 中的关联通过一个无限 Lorentz 变换与式~(\ref{Eq:correlationm}) 中的关联相联系~\cite{Ji:2013dva}。我们的``新''部分子形式体系处理不含时关联子与 IMF 波函数。由于算符不含时，只要在时间平移与 Lorentz 推进之间建立类比，它就是部分子物理的 Schr\"odinger 绘景~\cite{Ji:2020ect}。

然而在 QCD 中，关联 $\tilde h(\lambda)$ 与 $h(\lambda)$ 并不相同。差异源于 UV 截断的存在。在物理的动量分布中，截断总须远大于强子动量。结果，允许部分子动量大于强子动量，即 $|x|$ 可以大于 1，而不违反任何物理定律。另一方面，标准 PDF 具有支撑集 $|x|\le1$，对应于小于强子动量的 UV 截断。得益于渐近自由，这两个不同的 UV 极限可以通过 QCD 微扰论彼此连接起来。这使得从式~(\ref{Eq:correlatione}) 的欧氏形式中提取由式~(\ref{Eq:correlationm}) 定义的 LF 部分子物理成为可能。

式~(\ref{Eq:correlatione}) 是 LaMET 展开的出发点：我们先在有限但很大的动量 $P^z\gg\Lambda_{\rm QCD}$ 下计算准 LF 关联函数。为使展开有效，原则上需要 $-\infty<z<\infty$ 的 ${\tilde h}(z, P^z)$，或在所有准 LF 距离 $\tilde \lambda=z P^z$ 处的 ${\tilde h}(\tilde \lambda, P^z)$。当然现实中由于有限体积，格点数据总会停在某个大的 $z(\tilde \lambda)$ 处，记作 $z_{\rm L}(\tilde \lambda_{\rm L})$。有限 $\tilde \lambda_{\rm L}$ 带来的问题将在后文处理。本节的讨论假定在大 $P^z$ 下 $\tilde h(\tilde \lambda, P^z)$ 在整个 $\tilde \lambda$ 范围 $[-\infty,\infty]$ 内已知。

利用上述准 LF 关联可直接作傅里叶变换
\begin{equation}
    {\tilde f}(y,P^z) =\int^\infty_{-\infty} \frac{d\tilde \lambda}{2\pi}\, e^{i\tilde \lambda y}\, \tilde h(z, P^z)\ .
\end{equation}
${\tilde f}(y,P^z)$ 的物理解释取决于大动量展开~\cite{Ji:2013dva,Xiong:2013bka,Ma:2014jla,Izubuchi:2018srq}
\begin{align}\label{matching}
    {\tilde f}(y, P^z) &= \int^1_{-1} dx\, C\Big(\frac y x, \frac{xP^z}{\mu}\Big) f(x, \mu) \nonumber \\
    &+ {\cal O}\Big(\frac{\Lambda^2_{\rm QCD}}{y^2(P^z)^2}, \frac{\Lambda^2_{\rm QCD}}{(1-y)^2(P^z)^2}\Big),
\end{align}
其中 $\mu$ 是因子化标度，$\Lambda_{\rm QCD}$ 是强子能标，$f(x,\mu)$ 是可由上式提取的标准 PDF。大标度 $(yP^z)^2$ 与 $((1-y)P^z)^2$ 分别对应主动夸克与旁观者动量。按照标准 EFT 方法论，凡未被禁戒的大标度都应允许进入展开，而且由于时空对称性，线性依赖在维数正规化中不出现。因此该展开的有效性依赖于展开参量 $\Lambda^2_{\rm QCD}/[y^2(P^z)^2]$ 与 $\Lambda^2_{\rm QCD}/[(1-y)^2(P^z)^2]$ 的微小性。

上式中的因子 $C$ 可以微扰地计算。在 $\alpha_s$ 领头阶，可将 $\tilde f(y, P^z)$ 认同于 $f(y,\mu)$（暂且忽略幂次修正），于是二者在 $y\to 0$ 与 $y\to1$ 处具有相同的渐近行为。超出领头阶后，微扰修正会改变这一点。为看清这一点，取如下简单形式的 $f(x, \mu)$ 为例
\begin{equation}
x^a (1-x)^b,
\end{equation}
其中 $a, b$ 分别控制 $x\to 0$ 与 $x\to 1$ 处的渐近行为。单圈微扰修正使渐近行为发生如下改变
\begin{align}
\delta\tilde f(y, P^z)\sim \alpha_s y^a \ln y \ \  {\rm as}\ \  y\to 0\,.
\end{align}
当在微扰论中重求和至所有阶时，给出形如 $\tilde f(y, P^z)\sim  y^{a+\gamma}$ 的幂律行为，其中 $\gamma$ 与定义 $\tilde f(y, P^z)$ 的算符的反常量纲相关。对 $b>0$ 的真实 PDF，当 $y\to 1$ 时也出现类似行为。这也可以从下文给出的坐标空间分析中看出。

因此，对给定的大 $P^z$，存在一段 $y$ 区域，在该区域内高阶修正与幂次修正都很小，这一区域可以转译为 PDF 的有效 $x$ 范围。于是可以系统地得到区间 $[x_{\rm min}, x_{\rm max}]$ 内的 PDF（当 $P^z\to \infty$ 时，$x_{\rm min}$ 趋于 0、$x_{\rm max}$ 趋于 1）。换言之，LaMET 展开提供了一种在部分子动量 $x$ 的某个区间内计算部分子分布的自然途径，类似于在有限能量下从实验数据提取部分子分布。
